\section{Count and Say}
\label{sec:str-count-and-say}

\problemstatement{The count-and-say sequence starts with \texttt{"1"}. Each
term is generated by reading the previous term as runs of identical digits,
saying each run as its count followed by the digit. Given $n$, output the
$n$-th term. For example the terms are \texttt{1, 11, 21, 1211, 111221}.}

\begin{patternbox}
Each term is a \textbf{run-length encoding} of the previous one: scan the
previous string, group consecutive equal digits, and emit ``count digit''
for each group. Repeat $n-1$ times starting from \texttt{"1"}.
\end{patternbox}

\subsection*{Run-length encode, repeatedly}

Start with \texttt{"1"}. To produce the next term, walk the current string
counting the length of each maximal run of identical characters; append the
run length followed by the character. Applying this $n-1$ times yields the
$n$-th term. Each step reads and describes the previous one -- the sequence
is literally ``say what you see.''

\begin{ideabox}[Say the Runs You See]
The transformation from one term to the next is exactly run-length encoding:
consecutive duplicates collapse into a count-plus-digit pair. Iterating this
description operation generates the sequence.
\end{ideabox}

\subsection*{Algorithm}

\begin{enumerate}
  \item \texttt{cur = "1"}.
  \item Repeat $n-1$ times: scan \texttt{cur}, and for each run of length
        $c$ of digit $d$ append \texttt{c} then \texttt{d} to a new string;
        set \texttt{cur} to it.
  \item Return \texttt{cur}.
\end{enumerate}

\subsection*{Dry run}

\dryrun{$n=4$.}

\begin{itemize}
  \item \texttt{"1"} $\to$ one 1 $\to$ \texttt{"11"} $\to$ two 1s $\to$
        \texttt{"21"} $\to$ one 2 one 1 $\to$ \texttt{"1211"}.
  \item The $4$th term is \texttt{"1211"}.
\end{itemize}

\subsection*{C++ implementation}

\begin{lstlisting}
#include <bits/stdc++.h>
using namespace std;

string countAndSay(int n) {
    string cur = "1";
    for (int t = 1; t < n; t++) {
        string next;
        int i = 0, m = cur.size();
        while (i < m) {
            int j = i;
            while (j < m && cur[j] == cur[i]) j++;    // run of equal digits
            next += to_string(j - i);                 // count
            next += cur[i];                           // digit
            i = j;
        }
        cur = next;
    }
    return cur;
}

int main() {
    cout << countAndSay(4) << "\n";                // 1211
    cout << countAndSay(5) << "\n";                // 111221
    return 0;
}
\end{lstlisting}

\begin{complexitybox}
\textbf{Time Complexity.} $O(n\cdot L)$ where $L$ is the length of the
longest term (which grows roughly geometrically). \textbf{Space
Complexity.} $O(L)$.
\end{complexitybox}

\begin{takeawaybox}
Count-and-say is iterated run-length encoding -- each term describes the
runs of the previous one. Recognising the ``say what you see'' rule as RLE
turns a puzzling sequence into a simple grouping scan.
\end{takeawaybox}
